\subsection{Iterative Pruning Procedure}
\label{sec:iterative-method}

\paragraph{Algorithm overview.}
Algorithm~\ref{alg:iterative-pruning} specifies the iterative procedure
with input-independent gates. As in
Algorithm~\ref{alg:accuracy-guided-pruning}, $T_s,T_r$ are the search
and recovery budgets. The procedure repeats these phases for at most
$N$ cycles, carrying a set $D$ of permanently disabled channels.
The initialization policy $\mathcal I$ specifies the starting weights
for each phase. Numerical training and controller settings are given in
Sections~\ref{sec:mbv2-iterative-implementation}--\ref{sec:resnet-iterative-implementation}.

Let $(\theta_t,\phi_t)$ denote the saved search checkpoint, with
$\theta_t$ including weights and BatchNorm statistics and $\phi_t$ the gate
logits. Each cycle selects channels from the final checkpoint
$(\theta_{T_s},\phi_{T_s})$. For input-independent gates, the channel
scores satisfy $\overline r_{l,c}=r_{l,c}$.
The operator $\operatorname{Candidates}_{\geq1}(\phi_{T_s},D)$ considers
remaining channels with $r_{l,c}<\varepsilon_{\mathrm{prune}}$ in threshold
mode, or all remaining channels in budget-only mode. If every remaining
channel of a group qualifies, it excludes one with maximal score,
retaining at least one channel per group. The experimental tables denote
$\varepsilon_{\mathrm{prune}}$ by $\tau$.

Selection processes candidates in increasing score order. The optional
budget is $B=\rho P_{\mathrm{probe}}$, where $P_{\mathrm{probe}}$ counts
parameters in the full-width gated model. The estimated removal cost
$c(l,c)$ includes associated convolutional and normalization parameters.
These additive costs are approximate; $\rho$ is a removal cap rather
than an achieved compression ratio. Candidates exceeding the remaining
budget are skipped.

After selection, $\operatorname{Compact}(\theta_{T_s},m)$ physically
removes discarded channels and gates using the cumulative mask $m$.
The operator $\operatorname{Initialize}(g,\mathcal I)$ then replaces
its weights according to $\mathcal I$. Recovery uses a fresh optimizer
and scheduler for $T_r$ epochs and saves the checkpoint with the highest
validation accuracy. Only $D$ passes to the next cycle.

\begin{algorithm*}[t]
\caption[Iterative accuracy-guided structural channel pruning]{Iterative
accuracy-guided structural channel pruning. Each cycle performs search,
channel selection and recovery. Gate initialization, warmup, averaging
window and sparsity updates follow the configured search controller.
Setting $N=1$ gives the one-pass variant with the same selection rule.}
\label{alg:iterative-pruning}
\begin{algorithmic}[1]
\REQUIRE Network $f_\theta$; training and validation data;
dense reference accuracies $\{A_t^{\rm ref}\}_{t=1}^{T_s}$; initialization policy $\mathcal I$.
\REQUIRE Search and recovery budgets $T_s,T_r$; maximum cycles $N$;
accuracy tolerances $\delta_s,\delta_h$; selection mode, threshold
$\varepsilon_{\mathrm{prune}}$ and optional cap $\rho$.
\STATE $D\gets\varnothing$; recovery trajectory $\mathcal G\gets[\,]$.
\FOR{$k=1,\ldots,N$}
\STATE \textit{// Search}
\STATE Initialize $\theta$ from $\mathcal I$ and gate logits $\phi$;
permanently mask channels in $D$.
\STATE $\lambda_1\gets\lambda_0$; reset the search controller, optimizer and scheduler.
\FOR{$t=1,\ldots,T_s$}
\STATE Train $(\theta,\phi)$ for one epoch with $\lambda_t$;
measure validation accuracy $A_t$; save $(\theta_t,\phi_t)$.
\IF{$t\leq W$ or $t<w$}
\STATE $\lambda_{t+1}\gets\lambda_t$.
\ELSE
\STATE Update $\lambda_{t+1}$ using the configured validation-gap controller.
\ENDIF
\ENDFOR
\STATE \textit{// Channel selection}
\STATE $Q\gets\operatorname{Candidates}_{\geq1}(\phi_{T_s},D)$;
sort $Q$ by increasing $r_{l,c}$.
\STATE $R\gets\varnothing$;
$B\gets\rho P_{\mathrm{probe}}$ if a cap is set, otherwise $B\gets\infty$.
\FOR{channel $(l,c)$ in $Q$}
\IF{$c(l,c)\leq B$}
\STATE $R\gets R\cup\{(l,c)\}$; $B\gets B-c(l,c)$.
\ENDIF
\ENDFOR
\STATE $D\gets D\cup R$; $m_{l,c}\gets\mathbf 1[(l,c)\notin D]$.
\STATE \textit{// Recovery}
\STATE $g\gets\operatorname{Compact}(\theta_{T_s},m)$.
\STATE $g\gets\operatorname{Initialize}(g,\mathcal I)$;
initialize a fresh optimizer and scheduler.
\STATE Train $g$ for $T_r$ epochs without gates.
\STATE Append the best validation checkpoint of $g$ during recovery to $\mathcal G$.
\IF{$R=\varnothing$}
\STATE \textbf{break}.
\ENDIF
\ENDFOR
\RETURN Recovery trajectory $\mathcal G$; its last element is the final model.
\end{algorithmic}
\end{algorithm*}

For MobileNetV2, $\mathcal I$ reloads ImageNet weights for every phase,
using the surviving channel slices for recovery and a fresh classifier.
For iterative ResNet-50, every search and recovery phase starts with
random weights. Thus only $D$ passes between phases and cycles; recovery
does not inherit search weights or initialize the next search.
Candidates that exceed the remaining budget are skipped. The procedure
stops after recovery if no new removal was accepted, or after $N$ cycles.
Validation selects recovery checkpoints, rather than vetoing individual
channel removals. One-pass MobileNetV2 sets $N=1$ and retains these
selection and initialization rules.

